\documentclass[11pt]{letter}

\usepackage[margin=0.5in]{geometry}
\usepackage{hyperref}

\begin{document}

\begin{letter}{Scotiabank}

\opening{Dear Hiring Team,}

I am writing to apply for the Data Scientist, Fraud Threat Management position, Requisition ID 273753, with Scotiabank in Toronto. With a Master of Science in Computer Science from the University of Calgary and hands-on experience in Python, statistical model evaluation, machine learning, large-scale data processing, SQL-backed systems, and visualization, I am excited by the opportunity to help turn complex data into actionable insight for cyber-fraud prevention and response.

My graduate research focused on building data-driven software for complex Earth observation datasets. Working with BigGeo, Mitacs, and the University of Calgary, I developed Python and C++ pipelines for data ingestion, preprocessing, storage, retrieval, and analysis. I automated SAR and optical imagery preprocessing with GDAL, Rasterio, SNAP, and external data APIs, and built DGGS-based workflows that transformed satellite imagery into structured tensor files for downstream analysis. This gave me practical experience with large, imperfect real-world data, data quality, reproducible workflows, and the details that determine whether an analysis is trustworthy.

I also bring direct machine learning and model evaluation experience. In my DEM super-resolution project, I trained PyTorch ResNet models, collected and processed terrain data using Google Earth Engine and Rasterio, and evaluated results with statistical and terrain-aware metrics such as RMSE, slope, TRI, TPI, and wavelet-based errors. I also built a Streamlit-based workflow for data collection and preprocessing and a Polyscope-based visualization tool to inspect model outputs. That combination of modeling, evaluation, visualization, and communication aligns closely with a role focused on identifying patterns, testing hypotheses, improving detection capabilities, and explaining findings clearly.

My backend experience adds another useful layer for analytics work that needs robust, maintainable code. As a back-end developer, I built Django REST APIs for NLP, text-to-speech, and speech-to-text service plans backed by PostgreSQL, implemented authentication flows, and worked with structured service and payment data. As a teaching assistant at the University of Calgary, I supported students in big data applications, data-at-scale workflows, SQL, and database design, which strengthened my ability to communicate technical ideas to people with different backgrounds.

I have not worked directly in cyber-fraud incident management, but the analytical core of the role is very appealing to me: finding signal in large datasets, building reusable code, defining useful metrics, improving early detection, and helping teams make fast, evidence-based decisions. I would bring Scotiabank a strong foundation in Python, SQL, machine learning, statistical evaluation, data pipelines, Git-based development, and careful technical communication, along with genuine interest in learning the fraud and threat-management domain responsibly.

Thank you for considering my application. I would welcome the opportunity to discuss how my data science, machine learning, visualization, and large-scale data-processing experience can contribute to Scotiabank's Fraud Threat Management team.

\closing{Sincerely,

Amir Mirzai Golpayegani

\href{mailto:amirgolpaa24@gmail.com}{amirgolpaa24@gmail.com}}

\end{letter}

\end{document}
